\documentclass[11pt]{article}
\usepackage[a4paper,margin=1in]{geometry}
\usepackage[T1]{fontenc}
\usepackage[utf8]{inputenc}
\usepackage{lmodern}
\usepackage{amsmath,amssymb}
\usepackage{booktabs}
\usepackage{xcolor}
\usepackage{hyperref}
\hypersetup{
  colorlinks=true,
  linkcolor=blue!50!black,
  urlcolor=blue!50!black,
  citecolor=blue!50!black
}

\newcommand{\y}{\mathbf{y}}
\newcommand{\uvec}{\mathbf{u}}
\newcommand{\m}{\mathbf{m}}
\newcommand{\eps}{\boldsymbol{\epsilon}}
\newcommand{\M}{\mathbf{M}}
\newcommand{\A}{\mathbf{A}}
\newcommand{\R}{\mathbf{R}}
\newcommand{\G}{\mathbf{G}}
\newcommand{\Hmat}{\mathbf{H}}

\title{Critical Review of Semi-Synthetic Cell-Like Single-Color Controls for Spectreasy}
\author{Prepared for Spectreasy method development}
\date{\today}

\begin{document}
\maketitle

\section*{Executive Summary}

It is possible to create synthetic cell-like single-color controls (SCCs) by
starting from real unstained cell events and adding simulated fluorophore
spectra. This is a valid and potentially useful idea, but it should not be
described as creating ``perfect SCCs.'' A better name is
\textbf{semi-synthetic cell-like SCCs}. These controls are perfect only relative
to the spectrum and variation model assumed by the simulator. They cannot
discover an unknown true spectrum by themselves.

The idea is strongest as an \textbf{augmentation and benchmarking tool}. It can
make bead-derived spectra look more cell-like by adding them onto real cell AF
backgrounds. It can test whether Spectreasy recovers known simulated abundance.
It can also create controlled tandem-degradation or spectral-variant scenarios.
However, it should not replace real measured SCCs unless the spectral prior is
already trusted.

\section{Core Synthetic SCC Model}

Let an unstained cell event be
\begin{equation}
  \uvec_i \in \mathbb{R}^{D},
\end{equation}
where \(D\) is the number of spectral detectors. Let the fluorophore spectrum
for marker \(j\) be
\begin{equation}
  \m_j \in \mathbb{R}^{D}.
\end{equation}
A semi-synthetic positive event is generated as
\begin{equation}
  \y_i^{\text{synthetic}}
  =
  \uvec_i
  +
  a_i \m_j
  +
  \eps_i,
\end{equation}
where \(a_i\) is the simulated fluorophore abundance and \(\eps_i\) is optional
extra detector noise.

The most important feature is that \(\uvec_i\) is real. It contains real cell
autofluorescence, scatter-associated background, detector baseline behavior, and
event-to-event biological variation. The simulator does not need to invent a
cell background from scratch; it paints a fluorophore signal on top of a real
unstained cell.

\section{Why This Could Help Spectreasy}

\subsection{Cell-Like SCCs From Bead-Based Spectra}

Many SCCs are bead-based. Beads can provide clean bright spectra, but they do
not have the same AF background as cells. A semi-synthetic SCC can combine the
advantages of both:
\begin{equation}
  \text{bead-derived fluorophore spectrum}
  +
  \text{real unstained cell background}.
\end{equation}
This could help test how a fluorophore behaves when it is embedded in real cell
AF, even if the spectral prior came from beads.

\subsection{Known Truth for Benchmarking}

Because \(a_i\) is simulated, the true abundance is known. After unmixing, one
can compare the recovered abundance \(\hat{a}_i\) to the truth:
\begin{equation}
  e_i = \hat{a}_i - a_i.
\end{equation}
Useful benchmark metrics include:
\begin{equation}
  \operatorname{RMSE}
  =
  \sqrt{
    \frac{1}{n}
    \sum_{i=1}^{n}
    \left(
      \hat{a}_i - a_i
    \right)^2
  },
\end{equation}
and
\begin{equation}
  \operatorname{bias}
  =
  \frac{1}{n}
  \sum_{i=1}^{n}
  \left(
    \hat{a}_i - a_i
  \right).
\end{equation}
This makes synthetic SCCs useful for method development and regression tests.

\subsection{Controlled Variant and Tandem-Degradation Scenarios}

For tandem dyes or unstable fluorophores, the simulator can assign each event a
variant spectrum:
\begin{equation}
  \y_i
  =
  \uvec_i
  +
  a_i \m_{j,v_i}
  +
  \eps_i.
\end{equation}
For a simple degradation model, define
\begin{equation}
  \m_{j,v_i}
  =
  (1-\eta_i)\m_{\text{intact}}
  +
  \eta_i\m_{\text{degraded}},
\end{equation}
where
\begin{equation}
  0 \leq \eta_i \leq 1
\end{equation}
is the event-specific degradation fraction. If \(\eta_i = 0\), the event is
fully intact. If \(\eta_i = 1\), the event follows the degraded spectrum.

This lets developers create SCCs where not every positive event has the same
idealized spectrum. Spectreasy could then learn or test spectral variants under
controlled conditions.

\section{The Main Caveat: Synthetic SCCs Are Circular Without a Trusted Prior}

The central limitation is simple:
\begin{equation}
  \text{synthetic SCCs cannot invent the true fluorophore spectrum}.
\end{equation}
The simulator needs \(\m_j\) from somewhere. If the same assumed spectrum is
used to create synthetic SCCs and then recovered from those SCCs, the procedure
mostly validates its own assumption.

The useful direction is:
\begin{equation}
  \text{real SCC}
  \rightarrow
  \text{learn spectrum}
  \rightarrow
  \text{synthetic cell-like SCC augmentation}.
\end{equation}
The circular direction is:
\begin{equation}
  \text{assumed spectrum}
  \rightarrow
  \text{synthetic SCC}
  \rightarrow
  \text{``validated'' spectrum}.
\end{equation}
The second path is not independent validation. It can still be useful for
software testing, but not for proving that the spectrum is biologically correct.

\section{Why Averaging Broad Variant Mixtures Can Be Harmful}

If synthetic positives contain a broad mixture of variants, one might be tempted
to average them into one reference row:
\begin{equation}
  \bar{\m}_j
  =
  \frac{1}{n}
  \sum_{i=1}^{n}
  \m_{j,v_i}.
\end{equation}
This can be misleading. The average spectrum may be a spectrum that no real
event actually has. If the method then uses this average row for all events,
individual events may be poorly fit even though the average looks reasonable.

A safer approach is:
\begin{enumerate}
  \item keep a clean base spectrum as the main reference row;
  \item generate synthetic variant-positive events;
  \item use those events to populate or test the spectral-variant library;
  \item during sample unmixing, allow per-cell matching to variant spectra.
\end{enumerate}

In symbols, the matrix keeps the base row
\begin{equation}
  \m_j
\end{equation}
while the variant library stores
\begin{equation}
  \left\{
    \m_{j,1},
    \m_{j,2},
    \ldots,
    \m_{j,V}
  \right\}.
\end{equation}
The event-level solver can then select or test variants without replacing the
base biological reference with an artificial average.

\section{Noise and Intensity Models}

\subsection{Abundance Model}

A realistic synthetic SCC should not use one fixed abundance for all events.
Common choices include a lognormal model:
\begin{equation}
  a_i \sim \operatorname{LogNormal}(\mu,\sigma^2),
\end{equation}
or a mixture model with negative and positive cells:
\begin{equation}
  z_i \sim \operatorname{Bernoulli}(\pi),
\end{equation}
\begin{equation}
  a_i =
  \begin{cases}
    0, & z_i = 0, \\
    \operatorname{LogNormal}(\mu,\sigma^2), & z_i = 1.
  \end{cases}
\end{equation}
Here, \(\pi\) is the synthetic positive fraction.

\subsection{Detector Noise}

The simplest version adds no extra noise:
\begin{equation}
  \eps_i = \mathbf{0}.
\end{equation}
This is a reasonable starting point because the unstained event already carries
real background variation. A more explicit noise model could be:
\begin{equation}
  \epsilon_{id}
  \sim
  \mathcal{N}
  \left(
    0,
    \sigma_{0d}^{2}
    +
    \sigma_{1d}
    \left|
      y_{id}
    \right|
  \right).
\end{equation}
Here \(\sigma_{0d}^{2}\) is detector-specific baseline variance and
\(\sigma_{1d}\) controls signal-dependent variance.

This should be optional. Extra noise can make the simulation look more realistic
but also adds assumptions that are hard to verify.

\section{Implementation Rules}

The simulator should follow strict rules:
\begin{itemize}
  \item use raw linear detector values, not transformed or asinh values;
  \item preserve FSC, SSC, Time, and other non-spectral parameters from the
        unstained cell event;
  \item add fluorophore signal only to spectral detector channels;
  \item write metadata declaring the synthetic origin and the spectrum source;
  \item keep the original simulated abundance \(a_i\) in a truth table, not
        necessarily inside the FCS file;
  \item never label synthetic controls as equivalent to measured SCCs.
\end{itemize}

The detector update is:
\begin{equation}
  \y_{i,\text{detectors}}^{\text{new}}
  =
  \uvec_{i,\text{detectors}}
  +
  a_i\m_j.
\end{equation}
All non-detector channels are copied:
\begin{equation}
  \y_{i,\text{scatter/time}}^{\text{new}}
  =
  \uvec_{i,\text{scatter/time}}.
\end{equation}

\section{Suggested Spectreasy API}

\begin{verbatim}
simulate_scc_controls(
  unstained_fcs,
  reference_matrix,
  fluorophores = NULL,
  abundance_model = "lognormal",
  abundance_median = 5000,
  abundance_cv = 0.5,
  positive_fraction = 0.8,
  variant_model = "none",
  variant_library = NULL,
  noise_model = "none",
  output_dir = "synthetic_scc"
)
\end{verbatim}

For variant simulation:
\begin{verbatim}
simulate_scc_controls(
  unstained_fcs,
  reference_matrix,
  variant_model = "tandem_degradation",
  degradation_fraction = c(0, 0.4)
)
\end{verbatim}

The output should include:
\begin{itemize}
  \item synthetic FCS files;
  \item a control mapping CSV;
  \item a truth table with event-level \(a_i\), fluorophore, variant ID, and
        degradation fraction;
  \item metadata recording the source reference matrix and simulator settings.
\end{itemize}

\section{Best Use Cases}

\begin{table}[h!]
\centering
\begin{tabular}{@{}lll@{}}
\toprule
Use case & Value & Caveat \\
\midrule
Benchmarking & High & Tests assumed truth only \\
Bead-to-cell-like SCC augmentation & High & Spectrum still comes from beads \\
Tandem degradation experiments & High & Requires realistic degradation model \\
Missing real SCC replacement & Medium to low & Circular without trusted prior \\
Unknown spectrum discovery & Low & Synthetic data cannot reveal unknown truth \\
\bottomrule
\end{tabular}
\end{table}

\section{Critical Recommendation}

This feature is worth implementing as a simulator, augmentation, and benchmark
tool. It should not be marketed as a way to always create perfect SCCs. The most
accurate framing is:
\begin{quote}
Semi-synthetic cell-like SCCs create realistic cell-background controls from a
trusted spectral prior. They are useful for testing, augmentation, and controlled
variant scenarios, but they do not replace measuring the true spectrum when the
true spectrum is unknown.
\end{quote}

For Spectreasy, the best integration would be:
\begin{verbatim}
scc_mode = "semi_synthetic_cell"
\end{verbatim}
or, more explicitly,
\begin{verbatim}
augment_scc_with_synthetic_cells = TRUE
\end{verbatim}
with a metadata flag:
\begin{verbatim}
reference_source = "synthetic_from_prior"
\end{verbatim}

\section*{Bottom Line}

Yes, semi-synthetic cell-like SCCs are possible and useful. They can make
controls more cell-like, support benchmarking, and generate controlled spectral
variant scenarios. But they are not automatically perfect SCCs. They are only
as true as the spectrum and variant model used to create them.

\end{document}
